\renewcommand{\chaptername}{Ap\'endice}
\titleformat{\chapter}[block]{\normalfont}{}{0pt}{\tmchaphead}
\tmcapitulo{tmNavy}{Respuestas, definiciones y el mapa para llegar a los retos de TortuMath.}{x}

\chapter{Respuestas a ejercicios seleccionados}
\markboth{Ap\'endice A: Respuestas}{}
\tmgancho{Intenta cada ejercicio \textbf{antes} de mirar la respuesta. Si no te sale, vuelve al ejemplo resuelto de la secci\'on.}

\subsection*{Cap\'itulo 1: Fundamentos matem\'aticos}
\textbf{1.1}\label{resp:1-1}\quad
(1) $A\cup B=\{1, 2, 3, 4, 5\}$, $A\cap B=\{3, 4\}$, $A\setminus B=\{1, 2\}$. (2) $\{x\in\mathbb Z^+\mid x\text{ par}\wedge x\le10\}$. (3) $2^5=32$. (4) $A^c=\{6, 7\}$; $A\,\Delta\,B=\{1, 2, 3, 6\}$. (5) $(A\cup B)^c=\{4\}$ y $A^c\cap B^c=\{3, 4\}\cap\{1, 4\}=\{4\}$. (6) $30+25-10=45$. (7) $120+90-40=170$ usan alguna; no usan ninguna: $30$. (8) El primero s\'i; el segundo no (no son disjuntos: comparten el 2).

\medskip\noindent\textbf{1.2}\label{resp:1-2}\quad
(1) No, s\'i (V), no (predicado abierto). (2) V, F, V, V. (3) ``Existe un usuario sin token v\'alido''. (4) Ambas columnas dan V, F, V, V. (5) $\neg p\vee q$. (6) Tautolog\'ia, contradicci\'on, contingencia. (7) Modus Tollens con $q\to r$ y $\neg r$ da $\neg q$; Modus Tollens con $p\to q$ y $\neg q$ da $\neg p$. (8) \texttt{if (!a || b)}.

\medskip\noindent\textbf{1.3}\label{resp:1-3}\quad
(1) 6. (2) S\'i. (3) No (falta $(2, 1)$). (4) $\begin{bmatrix}0&1&0\\0&0&1\\0&0&1\end{bmatrix}$. (5) Se agregan $(1, 3),(2, 4),(1, 4)$. (6) S\'i: en el Hasse, 1 abajo, 2 y 3 en medio, 6 arriba. (7) Reflexiva, sim\'etrica y transitiva; 3 clases ($\{0, 3, 6\}$, $\{1, 4, 7\}$, $\{2, 5, 8\}$). (8) Todo el cierre son unos: $\begin{bmatrix}1&1&1\\1&1&1\\1&1&1\end{bmatrix}$ (el ciclo $1\to2\to3\to1$).

\medskip\noindent\textbf{1.4}\label{resp:1-4}\quad
(1) 7. (2) 3, 4 y 2. (3) No (porque $f(2)=f(-2)$); en $[0,\infty)$ s\'i. (4) 9 y 25. (5) $f^{-1}(x)=\frac{x+6}{3}$. (6) $5n+3\le5n+3n=8n$ para $n\ge1$: $c=8$, $n_0=1$. (7) $\log n,\ n,\ n\log n,\ n^2,\ 2^n$. (8) $4n^2+10n+7\le21n^2$ para $n\ge1$: $c=21$, $n_0=1$.

\medskip\noindent\textbf{1.5}\label{resp:1-5}\quad
(1) 7. (2) $6xy$. (3) $(2, 4)$. (4) $x^2+y^2\le25$ (c\'irculo de radio 5). (5) $3x^2y^2-4$ y $2x^3y+5$. (6) Circunferencia de radio 2 centrada en el origen. (7) $(3{,}2;\ 1{,}6)$. (8) $0{,}5c+30>50\Rightarrow c>40$.

\medskip\noindent\textbf{Mini proyecto:} (1) $P(u)\equiv(u\in A)\wedge(V(u)=1)\wedge(R(c, i)\le50)$. (2) El bloqueo es $\neg P(u)\equiv(u\notin A)\vee(V(u)=0)\vee(R(c, i)>50)$; el bloqueo autom\'atico del servidor es simplemente $R(c, i)>50$. (3) $\partial R/\partial c=0{,}5$ y $\partial R/\partial i=10$: un intento fallido pesa 20 veces m\'as que un punto de CPU.

\medskip\noindent\textbf{Autoevaluaci\'on 1:}\label{resp:q1}\quad 1 b, 2 a, 3 b, 4 b, 5 c, 6 b, 7 c, 8 b, 9 b, 10 b.

\medskip
\subsection*{Cap\'itulo 2: Matem\'atica discreta}
\textbf{2.1}\label{resp:2-1}\quad
(1) $4\cdot3=12$. (2) $4\cdot3\cdot2=24$. (3) $10^6=1\,000\,000$. (4) $26^3\cdot10^3=17\,576\,000$. (5) $5+3=8$. (6) $62^8=218\,340\,105\,584\,896\approx2{,}18\times10^{14}$. (7) Safari: $1\cdot2=2$; los otros 3 navegadores: $3\cdot3\cdot2=18$. Total $20$. (8) Son 25 objetos en 24 cajas: por el palomar, alguna caja recibe 2 o m\'as.

\medskip\noindent\textbf{2.2}\label{resp:2-2}\quad
(1) $4!=24$. (2) $\binom{10}4=210$. (3) $\binom{12}5=792$. (4) $P(6, 2)=30$. (5) $P(8, 3)=336$ y $\binom83=56$. (6) $\binom{10}4=210$ y $P(10, 4)=5040$. (7) $\frac{5!}{2!}=60$ (la E se repite dos veces). (8) $\binom nr=\frac{n!}{r!(n-r)!}$ es sim\'etrica en $r$ y $n-r$. $\binom{10}3=120=\binom{10}7$, la diferencia es $0$.

\medskip\noindent\textbf{2.3}\label{resp:2-3}\quad
(1) $2, 6, 18, 54, 162, 486$. (2) $4n$. (3) $2^5=32$. (4) $140$. (5) $1, 1, 2, 3, 5, 8, 13, 21$. (6) $2^5-1=31$. (7) $T(n)=\frac{n(n+1)}2$; $T(4)=10$. (8) Porque la recurrencia da el orden de crecimiento (por ejemplo $n\log n$ frente a $n^2$) sin ejecutar el algoritmo.

\medskip\noindent\textbf{2.4}\label{resp:2-4}\quad
(1) Compuesto: $7\cdot13$. (2) $2^2\cdot3^2\cdot5$. (3) $84=2\cdot30+24$, $30=1\cdot24+6$, $24=4\cdot6$: $\operatorname{MCD}=6$. (4) $36$. (5) S\'i. (6) $252=2\cdot105+42$, $105=2\cdot42+21$, $42=2\cdot21$: $21$. (7) Si $n=ab$ con $a\le b$, entonces $a\le\sqrt n$; un divisor propio siempre aparece antes de $\sqrt n$. (8) Con $n=55$ cualquiera lo factoriza al instante; RSA real usa primos de cientos de d\'igitos.

\medskip\noindent\textbf{2.5}\label{resp:2-5}\quad
(1) 5. (2) 6. (3) 3. (4) $63\bmod5=3$. (5) $20\bmod16=4$. (6) 5. (7) $13$. (8) RSA usa exponenciaci\'on modular e inversos modulares; una blockchain tambi\'en requiere hash, firmas, consenso, redes y validaci\'on.

\medskip\noindent\textbf{2.6}\label{resp:2-6}\quad
(1) Todos tienen grado 2. (2) $14$. (3) Instalar $\to$ compilar $\to$ ejecutar pruebas. (4) $\begin{bmatrix}0&1&1&0\\1&0&0&1\\1&0&0&1\\0&1&1&0\end{bmatrix}$. (5) No. (6) $A\to E\to C\to B$, costo $6$. (7) $\binom82=28$. (8) En la red f\'isica las aristas son cables reales con costo y capacidad; en las relaciones entre usuarios son v\'inculos l\'ogicos (amistad, seguimiento).

\medskip\noindent\textbf{2.7}\label{resp:2-7}\quad
(1) Una (por ejemplo $B$--$D$). (2) $\binom42=6$ cables. (3) $2^6=64$. (4) Respuesta libre; por ejemplo, un anillo de $n$ servidores donde el siguiente de $i$ es $(i+1)\bmod n$.

\medskip\noindent\textbf{Actividades de cierre:} (1) $24$. (2) $35$. (3) $\Theta(n\log n)$, como MergeSort. (4) $21$. (5) $4$. (6) a (8): respuesta libre; verifica que tu modelo identifique grados, caminos y puntos cr\'iticos.

\medskip\noindent\textbf{Autoevaluaci\'on 2:}\label{resp:q2}\quad 1 b, 2 b, 3 a, 4 b, 5 b, 6 b, 7 a, 8 c, 9 c, 10 b.

\medskip
\subsection*{Cap\'itulo 3: \'Algebra}
\textbf{3.1}\label{resp:3-1}\quad
(1) $x=5$. (2) $x=5$. (3) $x=15$. (4) $x=4$ y $x=3$ (Vieta: $4+3=7$, $4\cdot3=12$). (5) $\Delta=-4<0$: ninguna ra\'iz real. (6) $x\in\{\pm1,\pm3\}$; $x=\frac{10}{3}$ o $x=-4$. (7) Desde 101 horas ($h>100$). (8) $5\times8$ cm.

\medskip\noindent\textbf{3.2}\label{resp:3-2}\quad
(1) $(5, 4)$. (2) $(3, 6)$. (3) $(7, 3)$. (4) Incompatible. (5) Compatible indeterminado: $(t+1, t)$. (6) $(3, 3, 4)$. (7) 100 b\'asicos y 100 premium. (8) $S_1=135$, $S_2=270$, $S_3=195$.

\medskip\noindent\textbf{3.3}\label{resp:3-3}\quad
(1) $A+B=\begin{bmatrix}1&1\\4&4\end{bmatrix}$, $3A=\begin{bmatrix}3&6\\9&12\end{bmatrix}$. (2) $A^T=\begin{bmatrix}1&3\\2&4\end{bmatrix}$. (3) Rango 1; no es invertible. (4) $AB=\begin{bmatrix}2&-1\\4&-3\end{bmatrix}$, $BA=\begin{bmatrix}-3&-4\\1&2\end{bmatrix}$; no son iguales. (5) $\det A=-2$ y $A^{-1}=\begin{bmatrix}-2&1\\\frac32&-\frac12\end{bmatrix}$. (6) $\begin{bmatrix}\frac15&\frac15\\-\frac35&\frac25\end{bmatrix}$. (7) $(3, 2)$. (8) Ambos lados dan $\begin{bmatrix}2&4\\-1&-3\end{bmatrix}$.

\medskip\noindent\textbf{3.4}\label{resp:3-4}\quad
(1) 13 y 1. (2) 40 y $\frac15$. (3) No: $\det=4-4=0$. (4) $13-4=9$. (5) $k=2$ o $k=-3$. (6) $x=2$, $y=3$. (7) OK $\to$ UA. (8) \'Area $5{,}5$; con $\det=3$ queda $16{,}5$.

\medskip\noindent\textbf{Mini proyecto:} $\det A=16$; $(x, y, z)=(5, 4, 3)$; $\det K=5$ y $K^{-1}\equiv\begin{bmatrix}6&15\\5&16\end{bmatrix}\pmod{26}$; NUBE $\to$ IPOR.

\medskip\noindent\textbf{Autoevaluaci\'on 3:}\label{resp:q3}\quad 1 b, 2 b, 3 c, 4 a, 5 b, 6 b, 7 c, 8 b, 9 c, 10 c.

\medskip
\subsection*{Cap\'itulo 4: C\'alculo}
\textbf{4.1}\label{resp:4-1}\quad
(1) 6. (2) 2. (3) 3. (4) Por la derecha $+\infty$, por la izquierda $-\infty$: no existe. (5) $\frac{(x+2)(x-1)}{x-1}=x+2\to3$. (6) $\frac52$. (7) 8; la tabla da $7{,}9$; $7{,}99$; $7{,}999$ por la izquierda y $8{,}001$; $8{,}01$; $8{,}1$ por la derecha. (8) $\frac{1}{\sqrt{x+4}+2}\to\frac14$.

\medskip\noindent\textbf{4.2}\label{resp:4-2}\quad
(1) $3$. (2) $3x^2-2$ en $x=2$ da $10$. (3) $f'(x)=e^x(1+x)$ y $f'(0)=1$. (4) $6(2x+1)^2$ en $x=1$ da $54$. (5) $y=3x-2$. (6) $\frac{x^2-2x-1}{(x-1)^2}$. (7) $v(3)=32$ m/s y $a=10$ m/s$^2$. (8) $\frac{2x}{x^2+1}$.

\medskip\noindent\textbf{4.3}\label{resp:4-3}\quad
(1) $f_x=2xy$, $f_y=x^2+3$. (2) $ye^{xy}$. (3) $\nabla f=(6, 8)$ y $\|\nabla f\|=10$. (4) $(0, 0)$ y $(1, 1)$; en $(0, 0)$ es silla ($D=-9$) y en $(1, 1)$ es m\'inimo ($D=27$, $f_{xx}=6$). (5) M\'inimo en $(2,-3)$ con valor $-13$. (6) Punto silla. (7) $-12$ y $-6$. (8) En la direcci\'on de $\nabla f=(y, x)=(3, 2)$.

\medskip\noindent\textbf{4.4}\label{resp:4-4}\quad
(1) $x=3$, m\'inimo con $f=-4$. (2) $x=2$, m\'aximo con $f=4$. (3) $x=4$, $C=4$. (4) $x=-1$ m\'aximo ($f=2$) y $x=1$ m\'inimo ($f=-2$). (5) Cuadrado de $15\times15$, \'area $225$ m$^2$. (6) $n=10$, $C=60$. (7) $x=10$, $G=100$. (8) Porque en cada paso se mueve contra el gradiente, que es la direcci\'on de mayor descenso, hasta llegar a un punto donde $\nabla J=0$.

\medskip\noindent\textbf{4.5}\label{resp:4-5}\quad
(1) $x^3+C$. (2) $2$. (3) $x^2+x+C$. (4) $12$. (5) $e-1$. (6) $2$. (7) $60$. (8) Riemann: $3{,}75$; exacto: $\frac83\approx2{,}667$.

\medskip\noindent\textbf{Mini proyecto:} (1) $8$ ms; sustituir da $0/0$ y hay que factorizar. (2) $x=4$ con $C(4)=4$ y $C''=2>0$. (3) $N(5)=60$.

\medskip\noindent\textbf{Autoevaluaci\'on 4:}\label{resp:q4}\quad 1 c, 2 b, 3 b, 4 b, 5 b, 6 b, 7 b, 8 b, 9 a, 10 c.

\medskip
\subsection*{Cap\'itulo 5: Probabilidad y estad\'istica}
\textbf{5.1}\label{resp:5-1}\quad
(1) $\frac12$. (2) $\frac14$. (3) $0{,}7$. (4) $0{,}5+0{,}4-0{,}2=0{,}7$. (5) $0{,}1^3=0{,}001$. (6) $\frac{0{,}2}{0{,}5}=0{,}4$. (7) $0{,}75$. (8) $\frac58\cdot\frac47=\frac{5}{14}\approx0{,}357$.

\medskip\noindent\textbf{5.2}\label{resp:5-2}\quad
(1) $3{,}5$. (2) $0{,}5$. (3) $1-0{,}4-0{,}35=0{,}25$. (4) $E[X]=0{,}7$ y $\operatorname{Var}(X)=0{,}61$. (5) $E[3X+2]=14$ y $\operatorname{Var}(3X+2)=9\cdot2=18$. (6) $7$ Bs. (7) $100\cdot0{,}01-2=-1$ Bs: se pierde 1 Bs en promedio. (8) $E[X^2]=\frac{91}6$ y $E[X]^2=\frac{49}4$; $\frac{182-147}{12}=\frac{35}{12}$.

\medskip\noindent\textbf{5.3}\label{resp:5-3}\quad
(1) $\binom42\cdot0{,}5^4=\frac6{16}=0{,}375$. (2) $np=3$. (3) Aproximadamente $0{,}68$. (4) $5\cdot0{,}1\cdot0{,}9^4=0{,}32805$. (5) $e^{-2}\approx0{,}1353$. (6) $z=2$. (7) $1-0{,}9^3=0{,}271$. (8) $\frac{e^{-4}\cdot16}{2}=8e^{-4}\approx0{,}1465$.

\medskip\noindent\textbf{5.4}\label{resp:5-4}\quad
(1) $8$. (2) $5$. (3) $3$. (4) $\bar x=4$, $s^2=\frac{8}{2}=4$ y $s=2$. (5) Rango $10$ y mediana $6$. (6) La mediana. (7) $\mathrm{IQR}=17-12=5$; las cercas son $[4{,}5;\ 24{,}5]$ y $43$ queda fuera. (8) La media aumenta en 5; la desviaci\'on no cambia.

\medskip\noindent\textbf{5.5}\label{resp:5-5}\quad
(1) Relaci\'on lineal positiva muy fuerte. (2) $\hat y=14$. (3) Al aumentar $x$, $y$ tiende a aumentar. (4) $\hat y=\frac13+1{,}5x$. (5) $50\pm1{,}96\cdot2=[46{,}08;\ 53{,}92]$. (6) $5{,}8$. (7) Respuesta libre; por ejemplo, el n\'umero de procesos y el consumo de CPU, ambos crecen con la carga. (8) No: $r$ no cambia al sumar una constante.

\medskip\noindent\textbf{Mini proyecto:} (1) $1-0{,}1^3=0{,}999$. (2) $e^{-3}\approx0{,}0498$. (3) $\approx0{,}9772$. (4) Media $17{,}5$; mediana $14{,}5$; el at\'ipico es $43$ ms.

\medskip\noindent\textbf{Autoevaluaci\'on 5:}\label{resp:q5}\quad 1 b, 2 b, 3 a, 4 b, 5 c, 6 a, 7 a, 8 b, 9 b, 10 b.


% =====================================================================
\tmcapitulo{tmNavy}{Veinte problemas que mezclan los cinco cap\'itulos, con respuestas al final.}{x}
\chapter{Banco de repaso integrador}
\markboth{Ap\'endice B: Banco de repaso}{}
\tmgancho{Los ex\'amenes y los trabajos reales no avisan en qu\'e cap\'itulo est\'a el problema. Practica reconociendo qu\'e herramienta usar.}

\begin{tmtip}Para cada problema, primero escribe \textbf{qu\'e herramienta} vas a usar (conteo, derivada, matriz\ldots) y luego resu\'elvelo. Reconocer el tipo de problema es la mitad del trabajo.\end{tmtip}

\section*{Fundamentos}
\begin{tmlista}
\item \tmbasico\ Una plataforma tiene 150 usuarios: 90 usan modo oscuro, 70 usan notificaciones y 30 usan ambas. \textquestiondown Cu\'antos no usan ninguna?
\item \tmmedio\ Niega: ``todo servidor activo tiene certificado v\'alido o est\'a en mantenimiento''.
\item \tmmedio\ Demuestra que $7n^2+3n+1=O(n^2)$ hallando $c$ y $n_0$.
\item \tmdificil\ Con $J=w_1^2+3w_2^2$, $w=(2, 1)$ y $\alpha=0{,}1$, calcula un paso de descenso del gradiente.
\end{tmlista}

\section*{Matem\'atica discreta}
\begin{tmlista}[start=5]
\item \tmbasico\ Una clave tiene 4 letras min\'usculas \textbf{distintas}. \textquestiondown Cu\'antas claves hay?
\item \tmmedio\ De 9 desarrolladores se eligen 3 y luego uno de ellos es l\'ider. \textquestiondown De cu\'antas formas?
\item \tmmedio\ Resuelve $T(n)=T(n-1)+2$ con $T(1)=3$ y calcula $T(10)$.
\item \tmdificil\ Un \'indice circular de 12 posiciones est\'a en la 9 y avanza 20. \textquestiondown D\'onde termina? Calcula adem\'as $\operatorname{MCD}(360, 84)$.
\end{tmlista}

\section*{\'Algebra}
\begin{tmlista}[start=9]
\item \tmbasico\ Un plan cobra Bs 40 fijos m\'as Bs 1,5 por GB. Si pagaste Bs 100, \textquestiondown cu\'antos GB usaste?
\item \tmmedio\ Resuelve $2x+3y=13$ y $x-y=-1$.
\item \tmmedio\ Halla $\det\begin{bmatrix}2&1\\5&3\end{bmatrix}$ y su inversa.
\item \tmdificil\ Con la clave de Hill $K=\begin{bmatrix}1&2\\3&5\end{bmatrix}$, comprueba que es v\'alida y cifra HI.
\end{tmlista}

\section*{C\'alculo}
\begin{tmlista}[start=13]
\item \tmbasico\ Calcula $\lim_{x\to2}\frac{x^3-8}{x-2}$.
\item \tmmedio\ Deriva $f(x)=x^2e^{-x}$.
\item \tmmedio\ Halla el m\'inimo de $f(x, y)=x^2+y^2-2x-4y+10$.
\item \tmdificil\ La tasa de peticiones es $r(t)=3t+1$. Calcula las acumuladas entre $t=0$ y $t=4$.
\end{tmlista}

\section*{Probabilidad y estad\'istica}
\begin{tmlista}[start=17]
\item \tmbasico\ Dos servidores independientes fallan con probabilidad $0{,}05$ y $0{,}1$. \textquestiondown Probabilidad de que al menos uno falle?
\item \tmmedio\ $X$ vale 1, 2, 3 con probabilidades $0{,}2$, $0{,}5$, $0{,}3$. Calcula $E[X]$.
\item \tmmedio\ Binomial con $n=8$ y $p=0{,}25$: halla $E[X]$ y $\operatorname{Var}(X)$.
\item \tmdificil\ Datos: $4, 8, 6, 5, 3, 10$. Calcula media, mediana y desviaci\'on est\'andar muestral.
\end{tmlista}

\begin{tmsistemas}
\textbf{Respuestas del banco.} (1) $150-(90+70-30)=20$. (2) Existe un servidor activo sin certificado v\'alido y que no est\'a en mantenimiento. (3) $7n^2+3n+1\le11n^2$ para $n\ge1$: $c=11$, $n_0=1$. (4) $\nabla J=(4, 6)$ y $w'=(1{,}6;\ 0{,}4)$. (5) $26\cdot25\cdot24\cdot23=358\,800$. (6) $\binom93\cdot3=252$. (7) $T(n)=2n+1$ y $T(10)=21$. (8) $29\bmod12=5$ y $\operatorname{MCD}=12$. (9) $40$ GB. (10) $(2, 3)$. (11) $\det=1$ e inversa $\begin{bmatrix}3&-1\\-5&2\end{bmatrix}$. (12) $\det K=-1\equiv25$, coprimo con 26; HI $=(7, 8)\to(23, 61)\equiv(23, 9)=$ XJ. (13) $12$. (14) $e^{-x}(2x-x^2)$. (15) M\'inimo en $(1, 2)$ con valor $5$. (16) $\int_0^4(3t+1)\,dt=24+4=28$. (17) $1-0{,}95\cdot0{,}9=0{,}145$. (18) $2{,}1$. (19) $E=2$ y $\operatorname{Var}=1{,}5$. (20) Media $6$, mediana $5{,}5$, $s=\sqrt{\frac{34}{5}}\approx2{,}61$.
\end{tmsistemas}

% =====================================================================
\tmcapitulo{tmNavy}{Un proyecto que usa las cinco ramas: tu entrega integradora.}{x}
\chapter{Proyecto final integrador}
\markboth{Ap\'endice C: Proyecto final}{}
\tmgancho{\textquestiondown Y si construyes, paso a paso, el sistema de monitoreo inteligente de un cl\'uster usando todo lo que aprendiste?}

\begin{tmdef}{Proyecto: sistema de monitoreo inteligente}
Dise\~na y analiza un cl\'uster de servidores. Cada etapa usa un cap\'itulo del libro y entrega un resultado verificable. Se hace en equipos de 3 o 4 personas.
\end{tmdef}

\begin{center}
\begin{tikzpicture}[every node/.style={font=\sffamily\scriptsize,rounded corners=4pt,inner sep=5pt,align=center,minimum width=2.1cm,text=white}]
  \node[fill=tmAzul] (a) at (0,0) {1\\Modelar};
  \node[fill=tmMorado] (b) at (2.6,0) {2\\Contar y\\conectar};
  \node[fill=tmVerde] (c) at (5.2,0) {3\\Recursos y\\cifrado};
  \node[fill=tmAmbar] (d) at (7.8,0) {4\\Optimizar};
  \node[fill=tmCoral] (e) at (10.4,0) {5\\Medir y\\predecir};
  \foreach \p/\q in {a/b,b/c,c/d,d/e}{\draw[-{Stealth},tmOro,line width=1.4pt] (\p) -- (\q);}
\end{tikzpicture}
\end{center}

\subsection{Las cinco etapas}
\begin{center}
\begin{tabular}{|c|p{4.4cm}|p{7.2cm}|}
\hline
\rowcolor{tmAzulClaro}\textbf{Etapa} & \textbf{Qu\'e haces} & \textbf{Entregable}\\ \hline
1 Fundamentos & Define los conjuntos de servidores (activos, en alerta, ca\'idos), las reglas de alerta con l\'ogica y la funci\'on de riesgo $R(c, i)$ & Reglas escritas con conectivos, simplificadas con De Morgan; gr\'afica de curvas de nivel del riesgo\\ \hline
2 Discreta & Modela la red como grafo, cuenta configuraciones de despliegue y halla rutas & Matriz de adyacencia, grados, camino m\'as corto (Dijkstra) y enlace cr\'itico\\ \hline
3 \'Algebra & Plantea el balance de recursos como $A\mathbf x=\mathbf b$ y cifra los registros (logs) con Hill & Soluci\'on del sistema con $\det A$, y un mensaje cifrado y descifrado\\ \hline
4 C\'alculo & Optimiza el n\'umero de servidores e integra la tasa de peticiones & Punto \'optimo con $C''$ y peticiones acumuladas en un intervalo\\ \hline
5 Probabilidad & Calcula la disponibilidad con r\'eplicas, modela llegadas con Poisson y analiza las latencias & Disponibilidad, $P(X=k)$, media, mediana, percentiles y una recta de regresi\'on\\ \hline
\end{tabular}
\end{center}

\subsection{Datos de partida sugeridos}
\begin{tmsistemas}
Cl\'uster de 8 servidores con $R(c, i)=0{,}5c+10i$; red como la de la secci\'on 2.7; recursos de 17 n\'ucleos, 23 GB y 21 centenas de GB de disco; costo $C(n)=2n+\frac{800}{n}$; tasa $r(t)=4t+2$; r\'eplicas con probabilidad de fallo $0{,}1$; latencias $N(200, 20^2)$ ms.
\end{tmsistemas}

\subsection{R\'ubrica}
\begin{center}
\begin{tabular}{|p{5.8cm}|c|p{5.2cm}|}
\hline
\rowcolor{tmNavy}\textcolor{white}{\textbf{Criterio}} & \textcolor{white}{\textbf{Peso}} & \textcolor{white}{\textbf{Qu\'e se valora}}\\ \hline
Correcci\'on matem\'atica & 35\% & Planteamiento y c\'alculos verificados\\ \hline
Aplicaci\'on a sistemas & 25\% & Interpretaci\'on en el contexto del cl\'uster\\ \hline
C\'odigo en Python & 15\% & Reproduce los resultados\\ \hline
Retos TortuMath & 10\% & Misiones completadas en los 4 laboratorios\\ \hline
Presentaci\'on (5 min) & 15\% & Claridad y uso de gr\'aficos\\ \hline
\end{tabular}
\end{center}

\begin{tmcuidado}Un resultado sin comprobaci\'on vale la mitad. En cada etapa incluye la verificaci\'on (por ejemplo, $A\mathbf x=\mathbf b$ con tu soluci\'on, o $C'(x)=0$).\end{tmcuidado}

\tmreto{4}{Antes de entregar, repasa los laboratorios: cada misi\'on valida un resultado de tu proyecto (l\'imite, optimizaci\'on e integral acumulada).}

% =====================================================================
\tmcapitulo{tmNavy}{Las herramientas de c\'odigo que usan los ejemplos del libro, en una sola gu\'ia.}{x}
\chapter{Python para este libro}
\markboth{Ap\'endice D: Python}{}
\tmgancho{No necesitas instalar nada: abre Google Colab, crea un cuaderno nuevo y pega los bloques de c\'odigo del libro.}

\section*{Preparar el entorno}
\begin{tmpasos}
\item Entra a \texttt{colab.research.google.com} con tu cuenta de Google y elige \textbf{Cuaderno nuevo}.
\item En una celda escribe \texttt{import numpy, sympy, scipy, matplotlib} y ejec\'utala. Si no da error, todo est\'a listo.
\item Copia un bloque del libro en una celda y presiona \textbf{Shift + Enter}. Compara la salida con el comentario del bloque.
\end{tmpasos}

\begin{tmtip}Si trabajas en tu computadora, instala todo con \texttt{pip install numpy sympy scipy matplotlib}. En Windows usa el s\'imbolo del sistema o PowerShell.\end{tmtip}

\section*{Qu\'e librer\'ia usar en cada cap\'itulo}
\begin{center}
\begin{tabular}{|c|p{3.6cm}|p{7.6cm}|}
\hline
\rowcolor{tmAzulClaro}\textbf{Cap.} & \textbf{Librer\'ia} & \textbf{Funciones \'utiles}\\ \hline
1 & Python puro, NumPy & \texttt{set}, operadores \texttt{|}, \texttt{\&}, \texttt{-}, \texttt{\^{}}; \texttt{np.array}\\ \hline
2 & \texttt{math}, \texttt{sympy} & \texttt{comb}, \texttt{perm}, \texttt{factorial}, \texttt{gcd}, \texttt{pow(a, b, m)}, \texttt{factorint}, \texttt{isprime}\\ \hline
3 & NumPy, SymPy & \texttt{np.linalg.solve}, \texttt{det}, \texttt{inv}; \texttt{Matrix(...).rref()}\\ \hline
4 & SymPy, Matplotlib & \texttt{limit}, \texttt{diff}, \texttt{integrate}, \texttt{solve}; \texttt{plt.plot}\\ \hline
5 & \texttt{statistics}, SciPy & \texttt{mean}, \texttt{median}, \texttt{stdev}; \texttt{stats.binom}, \texttt{poisson}, \texttt{norm}; \texttt{np.polyfit}\\ \hline
\end{tabular}
\end{center}

\section*{Tres programas completos}
\subsection{Camino m\'as corto con Dijkstra (cap\'itulo 2)}
\begin{tmcodigo}
import heapq

G = {"A": {"B": 7, "E": 1},
     "B": {"A": 7, "C": 3, "E": 8},
     "C": {"B": 3, "D": 6, "E": 2},
     "D": {"C": 6, "E": 7},
     "E": {"A": 1, "B": 8, "C": 2, "D": 7}}

def dijkstra(G, origen):
    dist = {v: float("inf") for v in G}
    dist[origen] = 0
    cola = [(0, origen)]
    while cola:
        d, u = heapq.heappop(cola)
        if d > dist[u]:
            continue
        for v, peso in G[u].items():
            if d + peso < dist[v]:
                dist[v] = d + peso
                heapq.heappush(cola, (dist[v], v))
    return dist

print(dijkstra(G, "A"))   # {'A': 0, 'B': 6, 'C': 3, 'D': 8, 'E': 1}
\end{tmcodigo}

\subsection{Simulaci\'on de disponibilidad: m\'etodo Monte Carlo (cap\'itulo 5)}
En lugar de calcular $1-0{,}1^3$, \textbf{simulamos} 200\,000 d\'ias y contamos en cu\'antos sigue activo al menos un servidor:

\begin{tmcodigo}
import random

random.seed(1)
N = 200_000
activos = 0
for _ in range(N):
    if any(random.random() > 0.1 for _ in range(3)):   # 3 replicas
        activos += 1

print(activos / N)    # ~ 0.999 (exacto: 0.999)
\end{tmcodigo}

\subsection{Resumen estad\'istico de latencias (cap\'itulo 5)}
\begin{tmcodigo}
import statistics as st
import numpy as np
from scipy import stats

lat = [10, 12, 12, 14, 15, 16, 18, 43]
print(st.mean(lat), st.median(lat), st.stdev(lat))   # 17.5 14.5 10.61
print(np.percentile(lat, 95))                        # p95
print(stats.norm.cdf(240, loc=200, scale=20))        # 0.9772
print(stats.poisson.pmf(0, mu=3))                    # 0.0498
\end{tmcodigo}

\begin{tmcuidado}Una simulaci\'on da un resultado \textbf{aproximado} que cambia con la semilla; la f\'ormula da el valor exacto. Usa la simulaci\'on para verificar tu c\'alculo, no para sustituirlo.\end{tmcuidado}

\begin{tmejercicios}[Practica con c\'odigo]
\begin{tmlista}
\item \tmbasico\ Verifica en Python que $P(10, 3)=720$ usando \texttt{math.perm}.
\item \tmmedio\ Escribe una funci\'on que calcule el m\'aximo com\'un divisor con el algoritmo de Euclides y comp\'aralo con \texttt{math.gcd}.
\item \tmmedio\ Cambia el grafo del Dijkstra para que el peso $E$--$D$ sea 1. \textquestiondown Cambia el camino $A\to D$?
\item \tmdificil\ Simula 200\,000 d\'ias con 2 r\'eplicas en vez de 3 y compara con el valor exacto $1-0{,}1^2$.
\end{tmlista}
\end{tmejercicios}

% =====================================================================
\tmcapitulo{tmNavy}{Los errores que m\'as se repiten, reunidos para repasar antes de un examen.}{x}
\chapter{Errores comunes}
\markboth{Ap\'endice E: Errores comunes}{}
\tmgancho{Los errores no son al azar: casi siempre son los mismos. Con esta lista puedes revisar tu propio trabajo en un minuto.}

\begin{center}
\renewcommand{\arraystretch}{1.35}
\begin{tabular}{|c|p{5.4cm}|p{6.2cm}|}
\hline
\rowcolor{tmCoral}\textcolor{white}{\textbf{Cap.}} & \textcolor{white}{\textbf{Error}} & \textcolor{white}{\textbf{C\'omo evitarlo}}\\ \hline
1 & Confundir $a\in A$ con $\{a\}\subseteq A$ & Pertenencia es de elemento a conjunto; inclusi\'on es entre conjuntos\\ \hline
1 & Pensar que $p\to q$ es falsa si $p$ es falsa & Solo es falsa con $p$ verdadera y $q$ falsa\\ \hline
1 & Creer que la rec\'iproca equivale a la implicaci\'on & Solo la contrapositiva es equivalente\\ \hline
1 & Usar Big-O como igualdad exacta & Es una cota superior; busca la m\'as ajustada\\ \hline
1 & Usar un $\alpha$ muy grande en el descenso del gradiente & Si diverge, reduce $\alpha$\\ \hline
2 & Multiplicar etapas que no son independientes & Cuenta solo combinaciones v\'alidas\\ \hline
2 & Usar $P(n, r)$ cuando hay repeticiones & Con repetici\'on son $n^r$\\ \hline
2 & Llamar primo al 1 & El 1 tiene un solo divisor positivo\\ \hline
2 & Usar Dijkstra con pesos negativos & Exige pesos no negativos\\ \hline
3 & No comprobar tras elevar al cuadrado & Verifica siempre: pueden aparecer soluciones extra\~nas\\ \hline
3 & Cambiar el signo de solo un t\'ermino al quitar par\'entesis & Cambian todos los t\'erminos de adentro\\ \hline
3 & Suponer $AB=BA$ & El producto de matrices no es conmutativo\\ \hline
3 & Usar una clave de Hill con $\det$ no coprimo con 26 & Verifica $\operatorname{MCD}(\det K, 26)=1$\\ \hline
4 & Olvidar la cadena al derivar & Multiplica por la derivada de lo de adentro\\ \hline
4 & Dar por m\'inimo todo punto con $f'=0$ & Clasifica con $f''$ (o con $D$ en dos variables)\\ \hline
4 & Omitir la constante $+C$ & En la integral indefinida siempre va\\ \hline
5 & Confundir $P(A\mid B)$ con $P(B\mid A)$ & Usa Bayes\\ \hline
5 & Esperar que $E[X]$ sea un resultado posible & Es un promedio a largo plazo\\ \hline
5 & Dividir entre $n$ la varianza muestral & Se divide entre $n-1$\\ \hline
5 & Concluir causalidad de una correlaci\'on & Correlaci\'on no implica causalidad\\ \hline
\end{tabular}
\end{center}

\begin{tmtip}Cuando termines un problema, rev\'isalo con esta pregunta: \textquestiondown qu\'e error de esta tabla podr\'ia haber cometido? Si no encuentras ninguno, probablemente est\'a bien.\end{tmtip}

\begin{tmquiz}[Mini repaso: \textquestiondown qu\'e error es?]
\tmpreg{Alguien dice que $\sqrt{x+7}=x+1$ tiene soluciones $x=2$ y $x=-3$.}{No comprob\'o la ra\'iz extra\~na}{Olvid\'o la cadena}{Us\'o $P(n, r)$}{Divid\'o entre $n$}
\tmpreg{Alguien calcula $(2x+1)^3$ derivado como $3(2x+1)^2$.}{Olvid\'o la constante $C$}{Olvid\'o la regla de la cadena}{Us\'o Bayes mal}{No clasific\'o con $f''$}
\tmpreg{Alguien dice: ``si el correo es spam contiene la palabra con 60\%, entonces un correo con la palabra es spam con 60\%''.}{Confundi\'o $P(A\mid B)$ con $P(B\mid A)$}{Us\'o conteo}{Olvid\'o $+C$}{Us\'o Dijkstra}
\end{tmquiz}
\textbf{Respuestas:} 1 a, 2 b, 3 a.



% =====================================================================
\tmcapitulo{tmNavy}{Los t\'erminos clave del libro, explicados en una l\'inea.}{x}
\chapter{Glosario}
\markboth{Ap\'endice F: Glosario}{}

\begin{multicols}{2}
\raggedright
\sffamily\small
\textbf{Antiderivada.} Funci\'on cuya derivada es la dada.

\textbf{Aritm\'etica modular.} Aritm\'etica de residuos al dividir entre un m\'odulo.

\textbf{Bayes.} Regla que actualiza una probabilidad con nueva evidencia.

\textbf{Big-O.} Cota superior del crecimiento del costo de un algoritmo.

\textbf{Cardinalidad.} N\'umero de elementos de un conjunto.

\textbf{Cierre transitivo.} Relaci\'on m\'inima transitiva que contiene a otra.

\textbf{Combinaci\'on.} Selecci\'on sin importar el orden.

\textbf{Cofactor.} Menor con signo $(-1)^{i+j}$; sirve para Laplace.

\textbf{Correlaci\'on.} Medida $r\in[-1, 1]$ de relaci\'on lineal.

\textbf{De Morgan.} Leyes que niegan conjunciones y disyunciones.

\textbf{Derivada.} Tasa de cambio instant\'anea; pendiente de la tangente.

\textbf{Determinante.} N\'umero que indica si una matriz es invertible.

\textbf{Discriminante.} $\Delta=b^2-4ac$; anticipa las ra\'ices.

\textbf{Digrafo.} Grafo dirigido.

\textbf{Distribuci\'on normal.} Campana definida por $\mu$ y $\sigma$.

\textbf{Divisibilidad.} $a=bk$ para alg\'un entero $k$.

\textbf{Esperanza.} Valor promedio de una variable aleatoria.

\textbf{Factorial.} $n!=n(n-1)\cdots1$.

\textbf{Funci\'on biyectiva.} Inyectiva y sobreyectiva a la vez.

\textbf{Gauss.} M\'etodo de eliminaci\'on por filas.

\textbf{Gradiente.} Vector de derivadas parciales; apunta al m\'aximo crecimiento.

\textbf{Grafo.} Par $(V, E)$ de v\'ertices y aristas.

\textbf{Hash.} Funci\'on que mapea datos a un valor de tama\~no fijo.

\textbf{Inversa (matriz).} $A^{-1}$ con $AA^{-1}=I$.

\textbf{Integral.} Acumulaci\'on; \'area bajo la curva.

\textbf{L\'imite.} Valor al que se acerca $f(x)$ cuando $x\to a$.

\textbf{MCD.} Mayor divisor com\'un de dos enteros.

\textbf{Matriz.} Tabla de $m$ filas y $n$ columnas.

\textbf{Mediana.} Valor central de datos ordenados.

\textbf{Permutaci\'on.} Selecci\'on donde el orden importa.

\textbf{Poisson.} Distribuci\'on de eventos por intervalo con tasa $\lambda$.

\textbf{Primo.} Entero $>1$ con solo dos divisores positivos.

\textbf{Proposici\'on.} Enunciado verdadero o falso.

\textbf{Rango.} N\'umero de filas no nulas tras escalonar.

\textbf{Recurrencia.} Definici\'on de un t\'ermino por los anteriores.

\textbf{Regresi\'on lineal.} Recta de m\'inimos cuadrados $\hat y=a+bx$.

\textbf{Relaci\'on de equivalencia.} Reflexiva, sim\'etrica y transitiva.

\textbf{Riemann.} Suma de rect\'angulos que aproxima una integral.

\textbf{Sucesi\'on.} Lista ordenada de valores que sigue una regla.

\textbf{Varianza.} Promedio de las desviaciones al cuadrado.
\end{multicols}

% =====================================================================
\tmcapitulo{tmNavy}{Escanea, resuelve y acumula TortuPuntos.}{x}
\chapter{Mapa TortuMath}
\markboth{Ap\'endice G: Mapa TortuMath}{}
\tmgancho{Cada cap\'itulo tiene su propio laboratorio interactivo. Escanea el QR o escribe la direcci\'on.}

\begin{center}
\renewcommand{\arraystretch}{1.4}
\begin{tabular}{|c|p{3.2cm}|p{5.4cm}|c|}
\hline
\rowcolor{tmNavy}\textcolor{white}{\textbf{Cap.}} & \textcolor{white}{\textbf{Tema}} & \textcolor{white}{\textbf{Enlace}} & \textcolor{white}{\textbf{QR}}\\ \hline
1 & Fundamentos matem\'aticos & \texttt{\scriptsize\tmurlbr{1}} & \qrcode[height=2cm]{\tmurl{1}}\\ \hline
2 & Matem\'atica discreta & \texttt{\scriptsize\tmurlbr{2}} & \qrcode[height=2cm]{\tmurl{2}}\\ \hline
3 & \'Algebra & \texttt{\scriptsize\tmurlbr{3}} & \qrcode[height=2cm]{\tmurl{3}}\\ \hline
4 & C\'alculo & \texttt{\scriptsize\tmurlbr{4}} & \qrcode[height=2cm]{\tmurl{4}}\\ \hline
5 & Probabilidad y estad\'istica & \texttt{\scriptsize\tmurlbr{5}}\newline{\footnotesize (pendiente de publicaci\'on)} & \qrcode[height=2cm]{\tmurl{5}}\\ \hline
\end{tabular}
\end{center}

\begin{tmtip}Tambi\'en puedes entrar desde la p\'agina principal de la plataforma: \texttt{https://isaccmealla.github.io/TortuMath/}\end{tmtip}

\begin{tmsistemas}
\textbf{Fuentes para seguir aprendiendo:} ``Matem\'aticas para Ciencias de Computaci\'on'' (Lehman, Leighton y Meyer, MIT OpenCourseWare); ``Matem\'aticas Discretas: una introducci\'on abierta'' (Levin); recursos de criptograf\'ia del NIST (\texttt{nist.gov/cryptography}).
\end{tmsistemas}
